% -----------------------------------------------------------------------------------------------------
\section{重夸克偶素产生：NRQCD 因子化的进展}\label{sec:front-quarkonium}

NRQCD 因子化\cite{BodwinBraatenLepage1995} $\dd\sigma(H)=\sum_n\dd\hat\sigma_n(\Lambda)\langle\mathcal O^H_n(\Lambda)\rangle$ 对衰变是全阶论证（\cref{thm:heavy-nrqcd-decay}，[A]），对产生是猜想（\cref{conj:heavy-nrqcd-production}，[C]）。问题有二：(i) 短程系数的红外发散在所有阶是否与长程矩阵元的相匹配；(ii) 长程矩阵元的普适性。

\begin{proposition}[NNLO 的红外极点与规范完备的八重态算符 \tagP]\label{prop:front-nrqcd-ir}
在 NNLO（两圈）阶（均据摘要）：(a) 大 $p_T$ 胶子碎裂中，NNLO 的红外极点不被通常的色八重态产生算符匹配；规范不变性与因子化要求带 Wilson 线（非阿贝尔相位）的算符，用这些“规范完备”的矩阵元，因子化在 NNLO 恢复，先到 $\order{v^2}$，随后对八重态到单态的跃迁到 $v$ 的所有阶\cite{NayakQiuSterman2005a,NayakQiuSterman2005b,NayakQiuSterman2006}；保留：超出 NNLO，“in the absence of special cancellations”，可能需要 NRQCD 速度展开之外的非定域算符。(b) 对色八重态 S 波对碎裂为任意轨道角动量色单态对的长程矩阵元的协变两圈计算发现，红外发散与 Wilson 线方向无关，正如普适性所要求；结果在某些方面与 NQS 不同（来源已查明），两者都与 NRQCD 因子化相容，而“the general principle that underlies this confirmation \dots\ has yet to be revealed”\cite{Bodwin:2019bpf}。（本命题只涉及部分子层次的匹配。）
\uses{def:heavy-nrqcd}
\end{proposition}

\paragraph{地位与缺口。}[P]（两圈）。“NQS 的规范完备 Wilson 线在 NNLO 不改变结果”是过度表述——它们在 NNLO 是必需的。

\begin{remark}[唯象张力与替代方案]\label{rem:front-quarkonium-alternatives}
在 $1/p_T$ 中领头与次领头幂次的共线（碎裂）因子化（单部分子碎裂与重夸克对碎裂；\cite{Kang:2011mg,Kang:2014tta,Kang:2014pya,Fleming:2012wy,Ma:2013yla,Ma:2014eja,Bodwin:2014gia}）是核心陈述\cref{prop:heavy-quarkonium-fragmentation}（[A]）；它依赖于高 $p_T$ 单举共线因子化，而 Collins 在谈到详细证明的普遍缺乏时，把 NQS 2005 年的论证列为“notable exception”\cite[第~14.7~节，第~570--571~页]{Collins2011}。LHCb 的 $\eta_c$ 数据与重夸克自旋对称：一方认为它“challenges NRQCD factorization”\cite{Butenschoen:2014dra}，另一方得到相容的拟合\cite{Han:2014jya,Zhang:2014ybe}——两方并列，不作裁决；极化之谜仍然存在（此处不评述）。软胶子因子化\cite{Ma:2017xno}（据摘要）不先按相对速度展开就处理软强子化效应，并声称速度收敛更好——这是带部分检验的提议，作为定理属于猜想。
\end{remark}
